% chapters/conclusions/future_work.tex
% Sección 7.2: Trabajo futuro (~0.5-1 página)
% ============================================================================
%
% GUÍA: Proponer líneas de continuación del proyecto. Ideas:
%
%   1. Mejoras en el pipeline OCR:
%      - Integración de modelos de IA más avanzados (LLMs para OCR,
%        reconocimiento de escritura manuscrita con redes neuronales).
%      - Pre-procesamiento de imagen mejorado (corrección de perspectiva,
%        eliminación de ruido, binarización adaptativa).
%      - Detección automática de la estructura del examen sin plantilla.
%
%   2. Corrección asistida por IA:
%      - Sugerencias automáticas de calificación basadas en respuestas
%        anteriores similares.
%      - Detección de plagio entre exámenes.
%
%   3. Mejoras de escalabilidad:
%      - Despliegue en Kubernetes (migración desde Docker Compose).
%      - Migración a Amazon S3 / Azure Blob Storage.
%      - Base de datos gestionada (RDS, Cloud SQL).
%
%   4. Nuevos plugins de autenticación:
%      - OAuth2 / OpenID Connect (integración con proveedores de identidad
%        institucionales como LDAP de la UAM).
%      - SAML para federación de identidades.
%
%   5. Funcionalidades adicionales:
%      - Estadísticas y analíticas avanzadas de calificaciones.
%      - Exportación a plataformas LMS (Moodle, Canvas).
%      - Aplicación móvil para corrección.
%      - Soporte para exámenes tipo test (corrección automática completa).
%
%   6. Mejoras de infraestructura:
%      - Pipeline de CI/CD completo (tests, linting, despliegue automático).
%      - Monitorización y alertas (Prometheus, Grafana).
%      - Logs centralizados (ELK stack).
%
%   NOTA: El trabajo futuro debe ser realista y justificado. No listar
%   cosas genéricas; cada propuesta debe derivarse de una limitación
%   concreta del sistema actual o de un requisito no abordado.
% ============================================================================

% TODO: Redactar el trabajo futuro.
